\section{Методология численного моделирования}

\subsection{Математическая модель}

\subsubsection{Уравнение Рейнольдса}

Для тонкого смазочного слоя несжимаемой жидкости~\cite{pinkus, hamrock}:
\begin{equation}
    \frac{\partial}{\partial x}\!\left(h^3 \frac{\partial p}{\partial x}\right) + \frac{\partial}{\partial z}\!\left(h^3 \frac{\partial p}{\partial z}\right) = 6\eta U \frac{\partial h}{\partial x}.
    \label{eq:reynolds_dim}
\end{equation}

В уравнении~(\ref{eq:reynolds_dim}) левая часть описывает перенос давления через вязкий слой: первое слагаемое отвечает за окружной перенос (в направлении вращения $x$), второе -- за осевой перенос (в направлении оси подшипника $z$). Оба члена содержат $h^3$, что отражает сильную зависимость расхода от толщины плёнки (закон Пуазейля). Правая часть представляет собой источник давления -- она отлична от нуля только в области переменного зазора ($\partial h / \partial x \neq 0$) и описывает генерацию давления клиновым зазором~\cite{hamrock}. Именно этот член обеспечивает несущую способность подшипника: в зоне сужения зазора ($\partial h / \partial x < 0$) давление возрастает.

где \begin{tabular}[t]{@{}l@{\;--\;}l}
$h$ & толщина смазочного слоя;\\
$p$ & давление;\\
$\eta$ & динамическая вязкость;\\
$U = \pi D n / 60$ & линейная скорость поверхности вала~\cite{pinkus}.
\end{tabular}

\subsubsection{Безразмерная форма}

Вводя безразмерные переменные $\varphi = x/R$, $Z = z/L$, $H = h/c$, $P = pc^2/(6\eta U R)$, уравнение~(\ref{eq:reynolds_dim}) приводится к виду~\cite{hamrock}:
\begin{equation}
    \frac{\partial}{\partial\varphi}\!\left(H^3 \frac{\partial P}{\partial\varphi}\right) + \alpha^2 \frac{\partial}{\partial Z}\!\left(H^3 \frac{\partial P}{\partial Z}\right) = \frac{\partial H}{\partial\varphi},
    \quad \alpha^2 = \left(\frac{D}{L}\frac{\Delta\varphi}{\Delta Z}\right)^{\!2}.
    \label{eq:reynolds_ndim}
\end{equation}

Параметр $\alpha^2$ определяет соотношение окружного и осевого переноса давления. При $\alpha^2 \gg 1$ (короткий подшипник, $L \ll D$) доминируют осевые утечки, при $\alpha^2 \ll 1$ (длинный подшипник, $L \gg D$) -- окружной перенос. Для подшипника конечной длины оба механизма переноса существенны, и решение уравнения~(\ref{eq:reynolds_ndim}) является двумерным~\cite{pinkus}.

\subsubsection{Граничные условия}

Решение уравнения~(\ref{eq:reynolds_ndim}) требует задания граничных условий на границах расчётной области~\cite{hamrock}:

\begin{itemize}
    \item Периодичность: $P(0, Z) = P(2\pi, Z)$. Данное условие следует из замкнутости окружности подшипника -- давление должно быть непрерывной функцией угловой координаты.
    \item Атмосферное давление на торцах: $P(\varphi, \pm 1) = 0$. На торцах подшипника ($z = \pm L/2$, т.е. $Z = \pm 1$) давление равно атмосферному (принятому за нуль), так как торцы открыты в окружающее пространство.
    \item Условие кавитации (модель Гюмбеля): $P \geq 0$. В расширяющейся части зазора формальное решение уравнения Рейнольдса даёт отрицательные давления. Однако жидкость не способна выдержать значительных растягивающих напряжений -- в зоне отрицательного давления происходит кавитация (образование газовых пузырьков)~\cite{zhang2016}. В модели Гюмбеля отрицательные давления просто обнуляются.
\end{itemize}

\subsubsection{Зазор с микрорельефом}

Безразмерный зазор между валом и втулкой определяется суммой базового зазора и вклада микроуглублений~\cite{tala2012}:
\begin{equation}
    H(\varphi, Z) = 1 + \varepsilon\cos\varphi + \textstyle\sum_k \Delta H_k(\varphi, Z),
    \label{eq:clearance}
\end{equation}
где \begin{tabular}[t]{@{}l@{\;--\;}l}
$\varepsilon = e/c$ & относительный эксцентриситет;\\
$\Delta H_k$ & вклад $k$-го углубления.
\end{tabular}

Первые два члена описывают базовый зазор гладкого подшипника: единица соответствует радиальному зазору $c$, а $\varepsilon\cos\varphi$ -- смещению вала на эксцентриситет $e$ от центра втулки. Минимальный зазор $H_{\min} = 1 - \varepsilon$ достигается при $\varphi = \pi$, максимальный $H_{\max} = 1 + \varepsilon$ -- при $\varphi = 0$~\cite{pinkus}.

Третий член -- суммарный вклад всех углублений. Каждое углубление вносит добавку $\Delta H_k > 0$ в пределах своей области и $\Delta H_k = 0$ за её пределами. Конкретный вид функции $\Delta H_k$ зависит от типа профиля (эллипсоидальный, цилиндрический, параболический и др.) и определяется в расчётной части для каждого варианта.

\subsection{Численный метод решения}

\subsubsection{Дискретизация}

Расчётная область $\varphi \in [0, 2\pi]$, $Z \in [-1, 1]$ разбивается на равномерную сетку из $N_\varphi \times N_Z$ узлов с шагами $\Delta\varphi = 2\pi / N_\varphi$ и $\Delta Z = 2 / N_Z$~\cite{hamrock}. В каждом внутреннем узле $(i, j)$ уравнение~(\ref{eq:reynolds_ndim}) аппроксимируется методом конечных разностей второго порядка:
\begin{equation}
    A_{i,j} P_{i,j+1} + B_{i,j} P_{i,j-1} + C_{i,j} P_{i+1,j} + D_{i,j} P_{i-1,j} - E_{i,j} P_{i,j} = F_{i,j},
\end{equation}
где коэффициенты определяются через значения безразмерного зазора $H$ на полуцелых узлах (средние значения между соседними узлами -- центральная интерполяция):

$A = H_{j+\frac12}^3$ -- коэффициент при давлении в узле справа (по~$\varphi$);

$B = H_{j-\frac12}^3$ -- коэффициент при давлении в узле слева;

$C = \alpha^2 H_{i+\frac12}^3$ -- коэффициент при давлении в узле сверху (по~$Z$);

$D = \alpha^2 H_{i-\frac12}^3$ -- коэффициент при давлении в узле снизу;

$E = A + B + C + D$ -- диагональный коэффициент;

$F = \Delta\varphi(H_{j+\frac12} - H_{j-\frac12})$ -- правая часть (источник давления).

Здесь $H_{j+\frac12} = (H_{i,j} + H_{i,j+1})/2$ -- значение зазора на грани между узлами $j$ и $j+1$, аналогично для остальных полуцелых индексов.

\subsubsection{Метод Гаусса-Зейделя с релаксацией}

Полученная система линейных алгебраических уравнений решается итерационным методом Гаусса-Зейделя с верхней релаксацией (метод SOR -- Successive Over-Relaxation)~\cite{hamrock}:
\begin{equation}
    P_{i,j}^{(n+1)} = (1-\omega)\,P_{i,j}^{(n)} + \frac{\omega}{E_{i,j}} \bigl(A P_{i,j+1}^{(n)} + B P_{i,j-1}^{(n+1)} + C P_{i+1,j}^{(n)} + D P_{i-1,j}^{(n+1)} - F_{i,j}\bigr),
\end{equation}
где $\omega$ -- параметр релаксации.

Обход сетки выполняется последовательно по строкам: при обновлении давления в узле $(i, j)$ используются уже обновлённые значения в узлах $(i-1, j)$ и $(i, j-1)$, а также ещё не обновлённые значения в узлах $(i+1, j)$ и $(i, j+1)$. Такой порядок обхода обеспечивает более быструю сходимость по сравнению с методом Якоби~\cite{hamrock}.

Параметр верхней релаксации $\omega > 1$ ускоряет сходимость итерационного процесса. При $\omega = 1$ метод вырождается в стандартный метод Гаусса-Зейделя. Оптимальное значение $\omega$ зависит от размера сетки и свойств задачи; для данной задачи принято $\omega = 1{,}5$~\cite{pinkus}.

После каждого обновления применяется условие кавитации:
\begin{equation}
    P_{i,j}^{(n+1)} = \max(P_{i,j}^{(n+1)},\; 0).
\end{equation}

Итерации продолжаются до выполнения критерия сходимости по относительной невязке:
\begin{equation}
    \delta = \frac{\displaystyle\sum_{i,j} |P_{i,j}^{(n+1)} - P_{i,j}^{(n)}|}{\displaystyle\sum_{i,j} |P_{i,j}^{(n+1)}| + 10^{-15}} < 10^{-5}.
\end{equation}

Числитель представляет суммарное абсолютное изменение давления за одну итерацию, знаменатель -- текущий уровень давления (слагаемое $10^{-15}$ предотвращает деление на нуль). При $\delta < 10^{-5}$ поле давления считается сходившимся. Максимальное число итераций ограничено 50000.

\subsubsection{Интегральные характеристики}

По найденному полю давления $P(\varphi, Z)$ вычисляются интегральные характеристики подшипника~\cite{hamrock}.

Нагрузочная способность $F$ определяется как равнодействующая сил давления. Давление создаёт две компоненты силы -- горизонтальную $F_r^*$ (вдоль линии центров) и вертикальную $F_t^*$ (перпендикулярно):
\begin{equation}
    F = \frac{p_{\text{м}} R L}{2} \sqrt{F_r^{*2} + F_t^{*2}},
\end{equation}
где $F_r^* = \iint P\cos\varphi\,d\varphi\,dZ$, \; $F_t^* = \iint P\sin\varphi\,d\varphi\,dZ$, \; $p_{\text{м}} = 6\eta U R / c^2$ -- масштаб давления. Интегрирование выполняется по всей расчётной области методом трапеций.

Коэффициент трения $\mu$ характеризует отношение силы вязкого трения к несущей нагрузке:
\begin{equation}
    f^* = \iint \left(\frac{1}{H} + 3H\frac{\partial P}{\partial\varphi}\right) d\varphi\,dZ, \quad
    \mu = \frac{f^* \eta U R L}{c F}.
\end{equation}

Первый член подынтегрального выражения ($1/H$) соответствует сдвиговому (куэттовскому) течению, второй ($3H\,\partial P/\partial\varphi$) -- течению, обусловленному градиентом давления (пуазейлевскому)~\cite{pinkus}.

Расход смазки $Q$ определяет объём жидкости, вытекающей через торцы подшипника в единицу времени:
\begin{equation}
    Q^* = \iint \left(H - \frac{1}{2} H^3 \frac{\partial P}{\partial\varphi}\right) d\varphi\,dZ, \quad
    Q = Q^* U c R.
\end{equation}

Расход характеризует интенсивность обновления смазки в зазоре и определяет условия теплоотвода. Увеличение расхода способствует снижению температуры смазочного слоя~\cite{hamrock}.

\subsection{Программная реализация}

Численное решение уравнения Рейнольдса реализовано на языке Python с использованием библиотеки NumPy для работы с массивами и матричных операций. Для ускорения вычислений на графическом процессоре (GPU) применяется библиотека CuPy, обеспечивающая совместимый с NumPy интерфейс при выполнении вычислений на CUDA-ядрах.

Программа состоит из трёх основных модулей. Модуль формирования зазора реализует построение поля $H(\varphi, Z)$ для 10 типов профилей углублений (эллипсоидальный, параболический, цилиндрический, сферическая шапка и др.). Модуль решения выполняет итерационное решение дискретизированного уравнения Рейнольдса методом Гаусса-Зейделя с верхней релаксацией. Модуль постобработки вычисляет интегральные характеристики ($F$, $\mu$, $Q$) и формирует графические материалы.

Верификация программы выполнена путём сопоставления результатов для гладкого подшипника с известными аналитическими решениями~\cite{pinkus}. Проведена проверка сходимости при измельчении сетки: переход от сетки $250 \times 250$ к сетке $500 \times 500$ узлов приводит к изменению нагрузочной способности менее чем на $0{,}1$\%, что свидетельствует о достаточности выбранной дискретизации.

Параметры расчёта: сетка $500 \times 500$ узлов, параметр релаксации $\omega = 1{,}5$, критерий сходимости $\delta < 10^{-5}$, максимальное число итераций -- 50000. Расчёт одного значения эксцентриситета занимает порядка 10-30~с в зависимости от наличия GPU-ускорения.
